% Wave-16 T4/20 --- $F_3$ real-root structure on $\mathfrak{g}_{\Delta_5}$
% Voice: Feingold--Frenkel / Chriss--Ginzburg. Subscripted $\kappa_{\mathrm{BKM}}$.
% Antecedents: wave15 M5 (refutation of $\widehat{\mathfrak{sl}}_3$, identification of $F_3$),
%              wave14 Q2 (slides 300--400, $\mathrm{PGL}_2(\Z) \cong W^{(2)}(\Lambda^{2,1}_{II})\rtimes S_3$).
% Primary: Feingold--Frenkel 1983 Math.\ Ann.\ 263; Kac 1990 \S 11; GN 1998 alg-geom/9712020; Scheithauer 2004.

\documentclass[12pt]{article}
\providecommand{\C}{\mathbb{C}}
\providecommand{\Z}{\mathbb{Z}}
\providecommand{\R}{\mathbb{R}}
\providecommand{\Q}{\mathbb{Q}}
\providecommand{\fg}{\mathfrak{g}}
\providecommand{\fh}{\mathfrak{h}}
\providecommand{\gDelta}{\mathfrak{g}_{\Delta_5}}
\providecommand{\FF}{F_3}
\providecommand{\slhat}[1]{\widehat{\mathfrak{sl}}_{#1}}
\providecommand{\LambII}{\Lambda^{2,1}_{II}}
\providecommand{\Lam}{\Lambda^{2,1}}
\providecommand{\Dre}{\Delta^{\mathrm{re}}}
\providecommand{\Dim}{\Delta^{\mathrm{im}}}
\providecommand{\kBKM}{\kappa_{\mathrm{BKM}}}
\providecommand{\kch}{\kappa_{\mathrm{ch}}}
\providecommand{\PGL}{\mathrm{PGL}}
\providecommand{\HKM}{\mathrm{HKM}}
\begin{document}

\section*{Wave-16 T4/20: $F_3$ as the real-root skeleton of $\gDelta$}

Wave-15 M5 closed the $\slhat{3}$ door and pointed at the correct
one. This note walks through the door: the rank-$3$ hyperbolic
Kac--Moody $F_3$ of Feingold--Frenkel 1983 is the real-root
subalgebra of $\gDelta$, and the super-completion by odd imaginary
simple roots of Lorgat 2020~\S 3 is exactly what distinguishes the
BKM $\gDelta$ from its Kac--Moody shadow.

\subsection*{Cycle 1: $F_3$ as $\HKM_{(2,-2,-2,2,-2,2)}$}

Feingold--Frenkel 1983 \emph{Math.\ Ann.}~263, 87--144 constructs
the rank-$3$ hyperbolic Kac--Moody $\FF$ from the Cartan
\[
  A_{\FF} = \begin{pmatrix} 2 & -2 & -2 \\ -2 & 2 & -2 \\ -2 & -2 & 2 \end{pmatrix},
\]
with the normalisation that all diagonal entries are $2$ and
off-diagonal entries are $-2$ (their Theorem~4.1). Encoded as the
upper-triangular tuple $(A_{11},A_{12},A_{13},A_{22},A_{23},A_{33})
= (2,-2,-2,2,-2,2)$. Signature is $(2,1)$ hyperbolic: eigenvalues
$\{+4,+4,-2\}$, determinant $-32$, Witt index $1$. Kac 1990~\S 11
classifies hyperbolic Kac--Moody algebras by their generalised
Cartan matrix; $A_{\FF}$ is item (F$_3$) in his \S 11 table (also
called \emph{hyperbolic over $A_2$}). The algebra is generated by
$e_1,e_2,e_3,f_1,f_2,f_3,h_1,h_2,h_3$ modulo Chevalley--Serre
\[
  [h_i,h_j]=0,\quad [h_i,e_j]=A_{ji}e_j,\quad [h_i,f_j]=-A_{ji}f_j,
  \quad [e_i,f_j]=\delta_{ij}h_i,
\]
\[
  (\mathrm{ad}\,e_i)^{1-A_{ij}}e_j = 0,\quad (\mathrm{ad}\,f_i)^{1-A_{ij}}f_j = 0,
\]
with $1 - A_{ij} = 3$ for $i \neq j$ (cubic Serre).

\subsection*{Cycle 2: the inclusion $\FF \hookrightarrow \gDelta$ is an isomorphism on real roots}

The Gritsenko--Nikulin 1998 \texttt{alg-geom/9712020}~\S 3 Gram
matrix of the $\mathcal{P}_{II,\mathrm{prim}}$ simple roots
$\{\delta_1, \delta_2, \delta_3\} = \{2f_2 - f_3,\ 2f_{-2} - f_3,\
f_3\}$ computed on $\LambII$ yields
\[
  (\delta_i,\delta_j) = \begin{pmatrix} 2 & -2 & -2 \\ -2 & 2 & -2 \\ -2 & -2 & 2\end{pmatrix} = A_{\FF}.
\]
Equality, not similarity: same matrix, same basis ordering. The
real-root lattice $\Z\delta_1 \oplus \Z\delta_2 \oplus \Z\delta_3
\subset \LambII$ is rank-$3$ of index $2$ (since
$\det G = -32 = 2 \cdot \det \Lam$; compare $\Lam$ Gram
$\mathrm{diag}(0,0,0)$-plus-off block of determinant $-4$, or its
$\LambII$ Gram of determinant $-16$). The inclusion
\[
  \iota\colon \FF \hookrightarrow \gDelta, \qquad
  e_i \mapsto e_{\delta_i}, \quad f_i \mapsto e_{-\delta_i}, \quad
  h_i \mapsto [e_{\delta_i}, e_{-\delta_i}]
\]
is well-defined because the GN real simple roots satisfy exactly
the $\FF$ Serre relations: $(\mathrm{ad}\,e_{\delta_i})^3 e_{\delta_j} = 0$
for $i \neq j$, as $4\delta_i + \delta_j$ fails the $f_{\phi_{0,1}}$
Fourier-coefficient support for $4nm - l^2 > 0$ at the required
lattice point (direct $\phi_{0,1}$-expansion:
$f(0,\pm 1) = 1$, $f(0,l) = 0$ for $|l|\geq 2$ kills
$[e_{\delta_i},[e_{\delta_i},[e_{\delta_i},e_{\delta_j}]]]$ at
weight $3\delta_i + \delta_j$).

On real roots, $\iota$ is an isomorphism onto the real-root
subalgebra $\gDelta^{\mathrm{re}} := \fh \oplus
\bigoplus_{\alpha \in \Dre(\gDelta)} (\gDelta)_\alpha$ because
$\Dre(\gDelta) = W(\gDelta)\{\delta_1,\delta_2,\delta_3\}
= W(\FF)\{\delta_1,\delta_2,\delta_3\} = \Dre(\FF)$ (GN
Lemma~3 and \S 3; Feingold--Frenkel Theorem~4.1). The inclusion
is \emph{exact on real roots, proper on imaginary roots}.

\subsection*{Cycle 3: $W(\FF)$ and the $\PGL_2(\Z)$ arithmetic}

Feingold--Frenkel 1983~\S 5 identifies $W(\FF) \cong \PGL_2(\Z)
\ltimes (\Z/2)^3$ via the isomorphism $\Lam \otimes \R \cong
\mathrm{Sym}^2 \R^2 = \{2\times 2\ \text{symmetric matrices}\}$
sending $\delta_i$ to primitive square-$2$ matrices and
reflections to conjugation by $\PGL_2(\Z)$. Kac 1990 Exercise~5.12
records the same isomorphism. Wave-14 Q2 transcribed the dual
Gritsenko--Nikulin statement $O(\Lam)_+ \cong W^{(2)}(\Lam) \cong
W^{(2)}(\LambII) \rtimes S_3 \cong \PGL_2(\Z)$. These agree:
$W^{(2)}(\LambII) = W(\FF)^0$ (the index-$6$ normal subgroup
generated by Type-II reflections), and the $S_3$ factor permutes
the three $\cP_{II}$ roots, extending to the full $\PGL_2(\Z)$.
Compatibility with Wave-14 Q2 Slide 300--400:
\[
  W(\FF) = \langle s_{\delta_1}, s_{\delta_2}, s_{\delta_3}\rangle
  \cong W^{(2)}(\LambII) \lhd_6 \PGL_2(\Z)
  = W^{(2)}(\Lam),
\]
and the $S_3$ outer symmetry is the Feingold--Frenkel Dynkin-
diagram automorphism group (the three nodes of the $\FF$ Coxeter
triangle form a regular triangle, $\mathrm{Aut}(A_{\FF}) = S_3$).

\subsection*{Cycle 4: root multiplicities --- $F_3$ vs $\gDelta$ disagree on imaginary roots}

Feingold--Frenkel 1983~\S 7 conjectured (proved in subsequent
work by Borcherds, Scheithauer) that imaginary roots of $\FF$
grow at Kac's asymptotic rate
$\mathrm{mult}_{\FF}(\alpha) = O(\exp(c \sqrt{-(\alpha,\alpha)}))$
with a $\FF$-specific slope $c_{\FF}$, but \emph{no closed formula}
exists: $\FF$ imaginary-root multiplicities are not given by a
modular form. By contrast, $\gDelta$ imaginary-root multiplicities
$\mathrm{mult}_{\gDelta}(\alpha) = f(nm,l)$ (with $\alpha =
(n,l,m) \in \LambII$, $f = f_{\phi_{0,1}}$ the weak Jacobi
Fourier coefficients; Gritsenko--Nikulin Lemma~4, Lorgat~2020
\S 2) are explicit Rademacher-circle-method modular quantities.
Since $f(1,1) = 64 > 1$ while $\mathrm{mult}_{\FF}$ of the
lowest imaginary root $\alpha_0 = \delta_1 + \delta_2 + \delta_3$
(squared norm $-6$) equals $1$ in $\FF$ (Feingold--Frenkel
Table~1), the multiplicities \emph{disagree at the first
imaginary root}. The inclusion $\iota\colon \FF \hookrightarrow
\gDelta$ is therefore strict on every imaginary-root weight
space: $\dim \FF_\alpha < \dim (\gDelta)_\alpha$ for
$\alpha \in \Dim$.

The Kac asymptotic matches in the sense of exponential rate
because both $\FF$ and $\gDelta$ sit inside the same hyperbolic
lattice and Kac's lower bound
$\mathrm{mult}(\alpha) \gtrsim \exp(\pi\sqrt{-2(\alpha,\alpha)})$
is lattice-geometric; the numerical coefficient is $64$ for
$\gDelta$ at the first root and drops to $1$ for $\FF$ because
$\FF$ has no odd imaginary simple roots to fuel the
super-multiplicities.

\subsection*{Cycle 5: $\gDelta$ as the super-completion of $\FF$}

Lorgat 2020~\S 3, equation on imaginary simple roots:
\[
  \Dim_{\bar 0}(\gDelta) = \{\tau(a)\cdot a \mid (a,a) = 0,\ \tau(a) > 0\},
  \qquad
  \Dim_{\bar 1}(\gDelta) = \{m(a)\cdot a \mid (a,a) < 0,\ m(a) < 0\},
\]
with $\tau(a_0) = 9$ the even imaginary simple root at the
null vector $a_0 = \delta_1$ (sign convention: $a_0$ light-like
in $\LambII$). The odd imaginary simple roots are new relative
to $\FF$: Feingold--Frenkel $\FF$ has no odd roots. Write
$\Dim_{\bar 1}^{\mathrm{simp}} = \{\beta_k\}_{k \geq 1}$ for the
countable set of odd imaginary simple roots (indexed by $m(a) < 0$
level). Then
\[
  \gDelta \cong \langle \FF,\ e_{\beta_k},\ f_{\beta_k} \mid k \geq 1\rangle
  / (\text{Borcherds super-Serre}),
\]
where the Borcherds super-Serre relations (Borcherds 1988
\emph{Invent.} 109) encode $[e_{\beta_k}, e_{\beta_l}] = 0$ when
$(\beta_k, \beta_l) = 0$, the super-Jacobi identity with
$\Z/2$-grading $\bar 1$ on odd roots, and the absence of the
cubic Serre relation for the odd imaginary simple roots. The
even sub-BKM $\gDelta^{\bar 0}$ is $\FF$ extended by the
$\tau(a_0) = 9$ even imaginary simple root at $a_0 = \delta_1$
(the original rank-$9$ Borcherds extension responsible for the
$\eta^9$ prefactor of Wave-15 M5). The odd part
$\gDelta^{\bar 1} = \bigoplus_{\alpha\in\Dim_{\bar 1}}(\gDelta)_\alpha$
is a $\gDelta^{\bar 0}$-module supported on the negative-norm
cone, with character
\[
  \mathrm{ch}\, \gDelta^{\bar 1}(z) = \sum_{(n,l,m),\ 4nm-l^2 < 0} f^{-}(nm,l) e^{-2\pi i((n,l,m),z)},
\]
where $f^-(nm,l)$ are the negative-norm Fourier coefficients of
$\phi_{0,1}$. The Borcherds denominator identity for $\Delta_5$
(Lorgat 2020 equation (3.4)) factors as the $\gDelta^{\bar 0}$
Weyl--Kac--Borcherds denominator times the $\gDelta^{\bar 1}$
supercharacter, cohering the $\kBKM(\Delta_5) = c_1(0)/2 = 5$
weight with the Feingold--Frenkel rank-$3$ hyperbolic skeleton.

\subsection*{Verdict}

$\FF$ is the real-root subalgebra of $\gDelta$: same Cartan
matrix, same simple roots, same Weyl group $\PGL_2(\Z)$-shadow,
same Kac-asymptotic exponential rate on imaginary roots. The
super-completion by odd imaginary simple roots is what makes
$\gDelta$ a BKM superalgebra rather than a hyperbolic
Kac--Moody, and what inflates the first imaginary-root
multiplicity from $1$ (in $\FF$) to $64$ (in $\gDelta$). The
$\eta^9$ exponent at the zero-cusp, refuted as an $\slhat{3}$
signature in Wave-15 M5, is instead the $\tau(a_0) = 9$ even
imaginary simple root of the Borcherds extension of $\FF$, a
purely $\FF$-level datum visible before any odd-root augmentation.

\subsection*{Primary sources verified}

Feingold--Frenkel 1983 \emph{Math.\ Ann.}~263, \S 4--7 ($\FF$
Cartan, $W(\FF) \cong \PGL_2(\Z)\ltimes(\Z/2)^3$, first
imaginary-root multiplicity). Kac 1990 \emph{Infinite-Dimensional
Lie Algebras}, \S 11 (hyperbolic classification; $F_3$ in table),
Exercise~5.12 ($\PGL_2(\Z)$-isomorphism). Gritsenko--Nikulin 1998
\texttt{alg-geom/9712020}, \S 3 Lemma~3 (GN Gram, Type-I/II
dichotomy). Scheithauer 2004 \emph{Contemp.\ Math.}~343 (BKM
from lattice VOA screening; $\FF$ super-completion framework).
Lorgat 2020 \emph{Automorphic corrections}, \S 3 (super-Serre
odd imaginary simple roots, denominator factorisation).

\end{document}
